\section*{Reader's orientation}
\phantomsection\addcontentsline{toc}{section}{Reader's orientation}
\begin{center}
\begin{tabularx}{0.96\textwidth}{@{}lY@{}}
\toprule
\textbf{Status class} & \textbf{Meaning}\\
\midrule
Definition & Introduced structure or notation.\\
Theorem, proposition, corollary, lemma & Proved in this manuscript under the printed hypotheses.\\
Classical input & Imported standard theorem, stated explicitly because later proofs depend on it.\\
Boundary & A limitation, counterexample, or failure of an implication that is part of the formal scope.\\
\bottomrule
\end{tabularx}
\end{center}

\noindent\textbf{Formal and computational companions.} The formal companion explains the Lean development, its build procedure, theorem mapping, and the hypotheses of each verified statement. It distinguishes the 35-row foundational contract from the correspondence for the 108-record mathematical inventory. The computational companion presents those stable records, dependencies, and executable interfaces. Both companions are collected in the back matter. The mathematics below retains the full scope of its stated classical inputs and written proofs.


\begin{readerroadmap}
The argument has four stages. Reading them in order makes the later examples easier to follow.
\begin{enumerate}
\item \textbf{Grammar and realization.}  Specify which compositions of operators are legal, then distinguish that abstract syntax from the operators that actually act on a chosen space.
\item \textbf{Generation.}  Choose one or more seeds and ask which part of the realization is generated by their operator orbits.  The minimal number of required seeds is the seed rank.
\item \textbf{Packet coordinates.}  Resolve the space by a complete family of idempotents.  Every vector and every operator then has exact packet coordinates.
\item \textbf{Transport.}  Study how those coordinates behave under products, graphs, morphisms, quotients, completion, and realization functors.
\end{enumerate}
The first two layers describe what can be generated.  The last two describe how the generated object is resolved and transported.
\end{readerroadmap}

\subsection*{The basic conceptual chain}
\ConceptChain{\(\begin{gathered}
\text{typed grammar}\longrightarrow\text{realized algebra}\longrightarrow\text{seed hull}\\[.45em]
\longrightarrow\text{seed rank}\longrightarrow\text{packet coordinates}\\[.45em]
\longrightarrow\text{transfer blocks}\longrightarrow\text{functorial realization}.
\end{gathered}\)}

Each object in this chain has a different role.  A grammar is a presentation of allowed operator words.  A realized algebra is the image of that grammar in an endomorphism algebra.  A seed hull is an orbit span inside one particular module.  A packet system is an exact coordinate decomposition and need not be orthogonal.  A realization functor is required to preserve a declared mathematical profile, not merely the vocabulary.

\subsection*{Four questions answered by Chapter I}
\begin{center}
\begin{tabularx}{0.96\textwidth}{@{}>{\raggedright\arraybackslash}p{0.20\textwidth}Y Y@{}}
\toprule
\textbf{Question} & \textbf{Formal object} & \textbf{Principal answer}\\
\midrule
What operations are legal? & Quiver with relations or a unital operator algebra & Typed realizations are modules over the quotient path algebra.\\
What can a finite seed family generate? & Evaluation map \(A^{\oplus r}\to M\) & The generated hull is a quotient of a finite free module.\\
How many seeds are required? & \(\srank_A(M)\) & For a finitely generated module, pass to \(M/\Jac(A)M\). When the quotient algebra is semisimple Artinian, compute its simple multiplicities.\\
How are vectors and operators resolved? & Complete idempotents \((P_\alpha)\) & \(X\cong\bigoplus_\alpha P_\alpha X\) and \(T\leftrightarrow(P_\beta TP_\alpha)\).\\
\bottomrule
\end{tabularx}
\end{center}

\subsection*{Notation map}
\begingroup\small
\begin{center}
\begin{longtable}{@{}p{0.19\textwidth}p{0.27\textwidth}p{0.48\textwidth}@{}}
\toprule
\textbf{Symbol} & \textbf{Name} & \textbf{Meaning}\\
\midrule
\endhead
\(\mathfrak G=(Q,I)\) & Typed grammar & A quiver with finite vertex set and a two-sided ideal of source and target compatible relations.\\
\(\Lambda_{\mathfrak G}\) & Grammar algebra & The quotient path algebra \(\bbk Q/I\).\\
\(\A_X\) & Realized algebra & The image of the grammar algebra in \(\End_{\bbk}(|X|)\).\\
\(S\subseteq M\) & Seed family & The vectors from which the operator orbit is generated.\\
\(\Hull_A(S)\) & Generated hull & The submodule \(AS\), or its declared topological closure.\\
\(\srank_A(M)\) & Seed rank & The smallest cardinality of a generating family. A finite value is the minimum number of seeds.\\
\(\Pcal=(P_\alpha)\) & Packet system & Pairwise annihilating idempotents summing to the identity.\\
\(X_\alpha\) & Packet space & The range \(P_\alpha X\).\\
\(\mathcal A_{\Pcal}\) & Analysis & The map \(x\mapsto(P_\alpha x)_\alpha\).\\
\(\mathcal S_{\Pcal}\) & Synthesis & The map \((x_\alpha)\mapsto\sum_\alpha x_\alpha\).\\
\(T_{\beta\alpha}\) & Transfer block & The source and target block \(P_\beta TP_\alpha:X_\alpha\to X_\beta\).\\
\(\Supp_{\Pcal}(x)\) & Packet support & The packet indices on which \(x\) has a nonzero coordinate.\\
\(\psrank\) & Rank for seeds confined to individual packets & The minimal number of generating seeds each supported in a single packet.\\
\bottomrule
\end{longtable}
\end{center}
\endgroup

\subsection*{Small examples to keep in view}
The following small systems recur throughout the explanatory boxes. The worked models and visual atlas later develop these ideas in complete calculations.

\begin{microexample}[title={R1: one operator with several primary blocks}]
Let \(T\) act on a finite dimensional vector space.  Regard \(X\) as a \(\bbk[t]\)-module by \(t\cdot x=Tx\).  If \(T\) has three Jordan blocks for the same eigenvalue, then three seeds are necessary and sufficient for that primary component.  The largest number of blocks over any irreducible polynomial is the seed rank.
\end{microexample}

\begin{microexample}[title={R2: a matrix algebra with multiplicity}]
Let \(A=M_2(\bbk)\) act diagonally on \(M=(\bbk^2)^{\oplus3}\).  One seed can contribute at most two independent multiplicity directions, because the regular module contains two copies of the simple module.  Hence
\[
\srank_A(M)=\left\lceil\frac{3}{2}\right\rceil=2.
\]
This is the simplest nontrivial instance of the semisimple multiplicity formula.
\end{microexample}

\begin{microexample}[title={R3: an exact but oblique packet split}]
On \(X=\bbk^2\), set
\[
P_1=\begin{pmatrix}1&-1\\0&0\end{pmatrix},
\qquad
P_2=\begin{pmatrix}0&1\\0&1\end{pmatrix}.
\]
Then \(P_i^2=P_i\), \(P_1P_2=P_2P_1=0\), and \(P_1+P_2=I\).  The ranges form a direct sum, but over \(\mathbb R\) the projectors are not orthogonal in the Euclidean inner product.  Exact reconstruction and numerical conditioning are therefore separate questions.
\end{microexample}

\begin{microexample}[title={R4: a directed packet system}]
Let \(X=X_1\oplus X_2\oplus X_3\).  Suppose the generating operators permit transfers \(1\to2\) and \(2\to3\), with no path entering packet \(1\).  Packet \(1\) is the unique source component.  If the corner action on \(X_1\) is cyclic and the propagation hypotheses of \cref{thm:source-component-rank} hold, one seed in \(X_1\) can generate all three packets.
\end{microexample}

\subsection*{Suggested reading paths}
\begin{itemize}
\item For the minimal algebraic foundation, read Sections 1 to 5, then Sections 7 and 8.
\item For finite dimensional representation theory, concentrate on Sections 4 to 6.
\item For signal, wavelet, or quantum applications, read Sections 7 and 8 and the results on tensor products carefully.
\item For graphs and applications involving residual transport, read Sections 8 to 10.
\item For later categorical claims, read Sections 10 and 11 together with the final semantic realization criterion.
\end{itemize}

\begin{warningbox}
The chapter does not claim that every useful decomposition is automatically discovered, that every packet split is numerically stable, or that exact coordinates solve a nonlinear problem.  Chapter I establishes the linear language and its exact transport laws.  Refinement, obstruction theory, nonlinear trees, arithmetic valuations, and continuum limits belong to later chapters.
\end{warningbox}

